\documentclass[12pt,a4paper]{article}

% =========================
% Pacchetti
% =========================
\usepackage[utf8]{inputenc}
\usepackage[T1]{fontenc}
\usepackage[italian]{babel}
\usepackage{geometry}
\usepackage{setspace}
\usepackage{graphicx}
\usepackage{booktabs}
\usepackage{array}
\usepackage{xcolor}
\usepackage{hyperref}
\usepackage{amsmath}
\usepackage{float}
\usepackage{caption}
\usepackage{longtable}

% =========================
% Impostazioni pagina
% =========================
\geometry{margin=2.5cm}
\onehalfspacing

\hypersetup{
    colorlinks=true,
    linkcolor=black,
    urlcolor=blue,
    citecolor=black
}

% =========================
% Titolo
% =========================
\title{\textbf{Sistema di Raccomandazione per la Collision Avoidance Satellitare}\\
\large Supporto Decisionale Spiegabile per Eventi di Conjunction Spaziale}

\author{Samuele Pagone}

\date{\today}

\begin{document}

\maketitle

\begin{abstract}
Questo progetto sviluppa un sistema di raccomandazione spiegabile per il supporto decisionale nella collision avoidance satellitare. Il sistema analizza eventi di conjunction, cioè situazioni in cui un satellite target o un oggetto spaziale principale può avvicinarsi pericolosamente a un oggetto secondario, come un detrito spaziale, un satellite non operativo o un altro spacecraft. A partire dai dati degli eventi, il progetto costruisce un dataset a livello evento, crea feature legate al rischio e alla decisione operativa, calcola un punteggio globale di rischio, assegna uno score alle possibili azioni e raccomanda una tra quattro azioni: \texttt{no\_action}, \texttt{monitor}, \texttt{small\_maneuver} e \texttt{major\_maneuver}. La versione attuale combina un modulo di Machine Learning per la predizione del rischio con un recommendation layer rule-based e spiegabile. Il modello predittivo stima \texttt{max\_risk\_estimate}, mentre il recommendation layer trasforma rischio, distanza, urgenza, costo sintetico e priorità missione in una raccomandazione operativa.
\end{abstract}

\tableofcontents
\newpage

% =========================
\section{Introduzione}
% =========================

Gli eventi di conjunction satellitare si verificano quando due oggetti spaziali sono previsti passare a una distanza ravvicinata tra loro. In un contesto operativo, uno dei due oggetti è solitamente il target, cioè il satellite o l'oggetto principale da proteggere, mentre l'altro è il chaser o oggetto secondario, che può essere un detrito spaziale, un satellite non operativo o un altro satellite attivo.

Quando viene rilevato un possibile evento di conjunction, gli operatori devono decidere se non fare nulla, continuare a monitorare l'evento oppure effettuare una manovra di collision avoidance. Questa decisione non è banale, perché una manovra non necessaria può consumare carburante, interrompere la pianificazione della missione e introdurre complessità operative. Al contrario, una risposta tardiva o insufficiente può aumentare il rischio di collisione.

Lo scopo di questo progetto è costruire un sistema di raccomandazione per supportare questa decisione. Il sistema non si limita a stimare il rischio, ma trasforma indicatori fisici e operativi in una raccomandazione finale.

% =========================
\section{Obiettivo del progetto}
% =========================

L'obiettivo del progetto è raccomandare una tra quattro possibili azioni operative per ogni evento di conjunction:

\begin{itemize}
    \item \texttt{no\_action}
    \item \texttt{monitor}
    \item \texttt{small\_maneuver}
    \item \texttt{major\_maneuver}
\end{itemize}

Il sistema è pensato come un recommendation system spiegabile. Questo significa che ogni raccomandazione deriva da feature costruite esplicitamente, da punteggi interpretabili e da una logica decisionale trasparente.

La raccomandazione è destinata agli operatori di missione o a un flusso di supporto decisionale. Non deve essere interpretata come un comando autonomo inviato direttamente al satellite.

% =========================
\section{Descrizione del dataset}
% =========================

Il dataset principale utilizzato nel progetto è:

\begin{verbatim}
dataset/data/kelvins_competition_data/train_data.csv
\end{verbatim}

Il dataset originale contiene:

\begin{itemize}
    \item 162.634 righe
    \item 103 colonne
\end{itemize}

Il dataset contiene più osservazioni temporali per lo stesso \texttt{event\_id}. Questo significa che un singolo evento di conjunction può comparire più volte, con diversi valori di \texttt{time\_to\_tca} e con informazioni aggiornate sul rischio o sullo stato dell'evento.

I principali gruppi di variabili sono:

\begin{itemize}
    \item \texttt{event\_id}: identifica un evento di conjunction.
    \item \texttt{time\_to\_tca}: indica il tempo rimanente al closest approach.
    \item \texttt{risk} e \texttt{max\_risk\_estimate}: rappresentano informazioni sul rischio di collisione in forma logaritmica.
    \item \texttt{miss\_distance}: rappresenta la distanza minima prevista tra i due oggetti.
    \item \texttt{relative\_speed}: rappresenta la velocità relativa tra i due oggetti.
    \item colonne con prefisso \texttt{t\_}: descrivono il target object.
    \item colonne con prefisso \texttt{c\_}: descrivono il chaser object o oggetto secondario.
    \item colonne di posizione e velocità relativa: descrivono la geometria e il moto relativo dell'incontro tra i due oggetti.
\end{itemize}

Ogni riga finale a livello evento rappresenta quindi un incontro tra un oggetto target e un oggetto secondario, non lo stato di un singolo satellite isolato.

\subsection{Nota sulle unità di misura}

I file disponibili nel progetto non definiscono esplicitamente tutte le unità di misura fisiche. Di conseguenza, variabili come \texttt{miss\_distance}, \texttt{relative\_speed}, \texttt{time\_to\_tca} e le componenti di posizione relativa vengono interpretate nelle unità originali del dataset.

La logica di raccomandazione si basa principalmente su valori normalizzati e confronti relativi, senza assumere unità fisiche specifiche.

% =========================
\section{Preprocessing a livello evento}
% =========================

Il dataset originale non può essere usato direttamente come se ogni riga corrispondesse a una raccomandazione, perché lo stesso \texttt{event\_id} può apparire più volte. Nella prima versione del progetto, l'obiettivo è produrre una sola raccomandazione per ogni evento.

Per questo motivo, il progetto crea un dataset a livello evento selezionando, per ogni \texttt{event\_id}, la riga con il valore più piccolo non negativo di \texttt{time\_to\_tca}.

Questa regola permette di mantenere l'osservazione più vicina possibile al momento critico, ma ancora prima o al massimo al closest approach. I valori negativi di \texttt{time\_to\_tca} vengono esclusi perché indicano osservazioni probabilmente successive al momento di massimo avvicinamento.

Il dataset a livello evento contiene:

\begin{itemize}
    \item 13.143 eventi
\end{itemize}

Questa trasformazione è importante perché il sistema deve rispondere alla domanda:

\begin{quote}
Dato un evento di conjunction, quale azione deve essere raccomandata?
\end{quote}

% =========================
\section{Feature engineering}
% =========================

Il progetto crea diverse feature ingegnerizzate per rendere la logica di raccomandazione interpretabile e utile dal punto di vista operativo.

\subsection{Feature di probabilità}

Le colonne originali \texttt{risk} e \texttt{max\_risk\_estimate} sembrano essere rappresentate in scala logaritmica in base 10. Questo significa che valori più vicini a zero indicano una probabilità di collisione più alta.

Vengono quindi create le seguenti feature:

\begin{verbatim}
collision_probability = 10 ** risk
max_collision_probability = 10 ** max_risk_estimate
\end{verbatim}

Queste feature rendono i valori di rischio più interpretabili, mantenendo comunque le colonne originali per tracciabilità.

\subsection{Feature operative di rischio}

Vengono create due feature operative:

\begin{itemize}
    \item \texttt{distance\_risk}
    \item \texttt{urgency}
\end{itemize}

\texttt{distance\_risk} è una normalizzazione invertita di \texttt{miss\_distance}. Poiché una distanza minima più bassa indica un avvicinamento più critico, valori minori di \texttt{miss\_distance} producono valori più alti di \texttt{distance\_risk}.

\texttt{urgency} è una normalizzazione invertita di \texttt{time\_to\_tca}. Poiché un valore più basso di \texttt{time\_to\_tca} significa che manca meno tempo al closest approach, valori minori producono una maggiore urgenza.

\subsection{Feature normalizzata di velocità}

La feature \texttt{relative\_speed\_norm} viene creata normalizzando \texttt{relative\_speed} nell'intervallo da 0 a 1.

Questo passaggio è necessario perché i valori originali della velocità relativa sono numericamente molto più grandi delle altre feature normalizzate. La normalizzazione permette alla velocità di contribuire allo score senza dominare le altre variabili.

\subsection{Feature normalizzata di rischio}

La feature \texttt{risk\_norm} converte il rischio originale in scala logaritmica in uno score compreso tra 0 e 1:

\begin{itemize}
    \item valori vicini a 0 indicano rischio più basso;
    \item valori vicini a 1 indicano rischio più alto.
\end{itemize}

Questo rende l'indicatore principale di rischio compatibile con le altre feature normalizzate.

\subsection{Feature decisionali sintetiche}

Il dataset originale non contiene valori espliciti per costo carburante, delta-v o priorità della missione. Tuttavia, questi concetti sono importanti per costruire un sistema di supporto decisionale più realistico.

Per questo motivo, il progetto crea due feature sintetiche:

\begin{itemize}
    \item \texttt{fuel\_cost}
    \item \texttt{mission\_priority}
\end{itemize}

La feature sintetica \texttt{fuel\_cost} è definita come:

\begin{verbatim}
fuel_cost = 0.6 * relative_speed_norm + 0.4 * urgency
\end{verbatim}

Questa feature rappresenta un proxy della complessità operativa della manovra. Eventi con alta velocità relativa e alta urgenza vengono considerati più costosi o complessi dal punto di vista operativo.

La feature sintetica \texttt{mission\_priority} viene derivata deterministicamente da \texttt{mission\_id}. Produce valori tra 0.5 e 1.0 per simulare diversi livelli di importanza della missione.

Sia \texttt{fuel\_cost} sia \texttt{mission\_priority} sono assunzioni progettuali e non misure operative reali.

% =========================
\section{Costruzione del risk score}
% =========================

Il \texttt{risk\_score} rappresenta la severità complessiva dell'evento di conjunction. Combina rischio fisico, distanza, urgenza, velocità relativa e priorità della missione.

La formula utilizzata è:

\[
\text{risk\_score} =
0.35 \cdot \text{risk\_norm}
+ 0.25 \cdot \text{distance\_risk}
+ 0.20 \cdot \text{urgency}
+ 0.10 \cdot \text{relative\_speed\_norm}
+ 0.10 \cdot \text{mission\_priority}
\]

La feature \texttt{fuel\_cost} non viene inclusa nel \texttt{risk\_score}. Il costo carburante non aumenta la severità fisica dell'evento, ma influenza la desiderabilità delle possibili azioni. Per questo viene usato successivamente nello scoring delle azioni.

I pesi sono stati progettati nel seguente modo:

\begin{itemize}
    \item \texttt{risk\_norm} ha il peso maggiore perché rappresenta direttamente il rischio di collisione.
    \item \texttt{distance\_risk} è importante perché avvicinamenti più stretti sono più pericolosi.
    \item \texttt{urgency} conta perché meno tempo al TCA implica maggiore pressione operativa.
    \item \texttt{relative\_speed\_norm} contribuisce alla severità dell'incontro.
    \item \texttt{mission\_priority} aggiunge l'importanza operativa della missione.
\end{itemize}

% =========================
\section{Action scoring e logica di raccomandazione}
% =========================

Dopo aver calcolato il \texttt{risk\_score}, il sistema assegna uno score a ciascuna possibile azione:

\begin{itemize}
    \item \texttt{no\_action\_score}
    \item \texttt{monitor\_score}
    \item \texttt{small\_maneuver\_score}
    \item \texttt{major\_maneuver\_score}
\end{itemize}

La logica attesa è la seguente:

\begin{itemize}
    \item \texttt{no\_action} è favorita quando il rischio è basso.
    \item \texttt{monitor} è favorita per eventi a rischio intermedio.
    \item \texttt{small\_maneuver} è favorita per casi a rischio medio-alto.
    \item \texttt{major\_maneuver} è riservata a casi più severi, soprattutto quando la priorità della missione è alta.
\end{itemize}

La raccomandazione finale, \texttt{recommended\_action}, è selezionata come l'azione con lo score più alto.

In caso di pareggio esatto tra due o più azioni, viene scelta l'azione meno aggressiva secondo questo ordine:

\begin{verbatim}
no_action -> monitor -> small_maneuver -> major_maneuver
\end{verbatim}

Questa regola conservativa evita manovre non necessarie quando due azioni risultano ugualmente adatte.

% =========================
\section{Modulo di Machine Learning per la predizione del rischio}
% =========================

Oltre al recommendation system spiegabile, il progetto include un modulo di Machine Learning per la predizione del rischio.

Il modulo si trova nel file:

\begin{verbatim}
src/modeling.py
\end{verbatim}

L'obiettivo del modello non è predire direttamente l'azione finale, perché il dataset non contiene decisioni reali degli operatori come etichette. Il modello viene invece usato per predire una stima numerica del rischio dell'evento.

Il task di Machine Learning è quindi un task di regressione:

\begin{quote}
Predire \texttt{max\_risk\_estimate} a partire dalle feature fisiche e orbitali dell'evento di conjunction.
\end{quote}

\subsection{Target del modello}

La variabile target scelta è:

\begin{verbatim}
max_risk_estimate
\end{verbatim}

Questa variabile rappresenta una stima del rischio massimo dell'evento. Essendo un valore numerico continuo, il problema viene trattato come un problema di regressione.

\subsection{Feature di input}

Il modello utilizza feature fisiche e orbitali che descrivono lo scenario tra target object e chaser object.

Le feature usate sono:

\begin{itemize}
    \item \texttt{time\_to\_tca}
    \item \texttt{miss\_distance}
    \item \texttt{relative\_speed}
    \item \texttt{relative\_position\_r}
    \item \texttt{relative\_position\_t}
    \item \texttt{relative\_position\_n}
    \item \texttt{relative\_velocity\_r}
    \item \texttt{relative\_velocity\_t}
    \item \texttt{relative\_velocity\_n}
    \item \texttt{mahalanobis\_distance}
    \item \texttt{mission\_id}
    \item \texttt{c\_object\_type}
\end{itemize}

Per evitare data leakage, non vengono usate come input colonne derivate dal rischio o colonne prodotte dal sistema di raccomandazione.

Sono quindi escluse colonne come:

\begin{itemize}
    \item \texttt{risk}
    \item \texttt{max\_risk\_estimate}
    \item \texttt{collision\_probability}
    \item \texttt{max\_collision\_probability}
    \item \texttt{risk\_norm}
    \item \texttt{risk\_score}
    \item \texttt{recommended\_action}
    \item gli action scores
\end{itemize}

Questa scelta è importante perché impedisce al modello di ricevere informazioni che contengono direttamente o indirettamente la risposta da predire.

\subsection{Modelli utilizzati}

Sono stati confrontati due modelli:

\begin{itemize}
    \item \texttt{DummyRegressor}
    \item \texttt{RandomForestRegressor}
\end{itemize}

Il \texttt{DummyRegressor} viene usato come baseline. Questo modello predice sempre un valore semplice, vicino alla media del target, e serve per verificare se il modello principale riesce davvero a imparare informazioni utili dai dati.

Il \texttt{RandomForestRegressor} è il modello principale. È stato scelto perché funziona bene con dati tabellari, gestisce relazioni non lineari ed è robusto per una prima versione del modulo predittivo.

\subsection{Metriche di valutazione}

I modelli sono stati valutati con tre metriche di regressione:

\begin{itemize}
    \item \textbf{MAE}: errore medio assoluto;
    \item \textbf{RMSE}: radice dell'errore quadratico medio;
    \item \textbf{R2}: quota di variabilità del target spiegata dal modello.
\end{itemize}

I risultati ottenuti sono:

\begin{table}[H]
\centering
\begin{tabular}{lrrr}
\toprule
Modello & MAE & RMSE & R2 \\
\midrule
\texttt{DummyRegressor} & 0.8350 & 1.0488 & -0.0002 \\
\texttt{RandomForestRegressor} & 0.3802 & 0.5213 & 0.7529 \\
\bottomrule
\end{tabular}
\caption{Valutazione dei modelli di regressione per la predizione di \texttt{max\_risk\_estimate}.}
\end{table}

\subsection{Interpretazione dei risultati}

Il \texttt{DummyRegressor} ha un valore di \texttt{R2} circa pari a zero. Questo indica che la baseline non spiega realmente la variabilità del target.

Il \texttt{RandomForestRegressor}, invece, ottiene:

\begin{verbatim}
R2 = 0.7529
\end{verbatim}

Questo significa che il modello riesce a spiegare circa il 75\% della variabilità di \texttt{max\_risk\_estimate} nel test set.

Inoltre, il modello Random Forest ottiene valori di MAE e RMSE molto più bassi rispetto alla baseline. Questo conferma che le feature fisiche e orbitali dello scenario contengono informazione utile per stimare il rischio massimo dell'evento.

\subsection{Ruolo del Machine Learning nel sistema complessivo}

Il modulo di Machine Learning non sostituisce il recommendation system. Il suo ruolo è stimare il rischio.

La struttura complessiva del progetto diventa quindi:

\begin{verbatim}
Feature dell'evento di conjunction
        ↓
Modulo ML di predizione del rischio
        ↓
Stima di max_risk_estimate
        ↓
Recommendation layer spiegabile
        ↓
Azione raccomandata
\end{verbatim}

La raccomandazione finale continua a essere prodotta dal recommendation layer rule-based, che combina rischio, distanza, urgenza, velocità, costo sintetico e priorità missione per selezionare l'azione più adatta.

% =========================
\section{Output generati}
% =========================

La pipeline genera i seguenti output:

\begin{itemize}
    \item \texttt{outputs/tables/event\_level\_data.csv}: dataset completo processato, con colonne originali, feature ingegnerizzate, score e raccomandazioni.
    \item \texttt{outputs/recommendations/recommendations.csv}: file pulito finale con le colonne principali della raccomandazione.
    \item \texttt{reports/recommendation\_summary.md}: report riassuntivo con distribuzione delle raccomandazioni, percentuali, score medi ed esempi.
    \item \texttt{reports/recommendation\_validation.md}: report di validazione descrittiva che controlla la coerenza delle raccomandazioni.
    \item \texttt{reports/model\_evaluation.md}: report di valutazione del modello di Machine Learning per la predizione di \texttt{max\_risk\_estimate}.
    \item \texttt{outputs/models/risk\_prediction\_model.joblib}: modello \texttt{RandomForestRegressor} addestrato e salvato.
    \item \texttt{outputs/figures/}: grafici statici di esempio per le diverse classi di raccomandazione.
\end{itemize}

La funzione di visualizzazione genera un grafico di esempio per ogni raccomandazione disponibile:

\begin{itemize}
    \item \texttt{example\_no\_action.png}
    \item \texttt{example\_monitor.png}
    \item \texttt{example\_small\_maneuver.png}
    \item \texttt{example\_major\_maneuver.png}
\end{itemize}

Queste figure sono supporti visivi. Rappresentano in modo statico la geometria relativa tra target e chaser. Non simulano propagazione orbitale, delta-v o dinamiche successive alla manovra.

% =========================
\section{Risultati}
% =========================

La distribuzione finale delle raccomandazioni è:

\begin{table}[H]
\centering
\begin{tabular}{lrr}
\toprule
Azione raccomandata & Conteggio & Percentuale \\
\midrule
\texttt{monitor} & 10927 & 83.1393\% \\
\texttt{major\_maneuver} & 1945 & 14.7988\% \\
\texttt{small\_maneuver} & 209 & 1.5902\% \\
\texttt{no\_action} & 62 & 0.4717\% \\
\bottomrule
\end{tabular}
\caption{Distribuzione finale delle raccomandazioni.}
\end{table}

La raccomandazione più frequente è \texttt{monitor}. Questo indica che la maggior parte degli eventi richiede osservazione continua, ma non una manovra immediata.

La raccomandazione \texttt{major\_maneuver} compare per un gruppo più piccolo ma significativo di casi ad alto rischio o alta priorità.

La raccomandazione \texttt{small\_maneuver} compare in un numero inferiore di casi intermedi, dove una manovra è utile ma una manovra maggiore non è l'opzione più adatta.

La raccomandazione \texttt{no\_action} compare solo per casi a rischio basso.

% =========================
\section{Validazione e interpretazione}
% =========================

Il validation report conferma che le raccomandazioni sono coerenti con la logica degli score.

I valori medi delle feature per raccomandazione sono riassunti nella tabella seguente.

\begin{table}[H]
\centering
\begin{tabular}{lrrrrrr}
\toprule
Azione & Risk & Fuel & Priority & Urgency & Dist. Risk & Speed Norm \\
\midrule
\texttt{major\_maneuver} & 0.7602 & 0.6984 & 0.7281 & 0.7893 & 0.8660 & 0.6378 \\
\texttt{monitor} & 0.5152 & 0.6905 & 0.6611 & 0.8222 & 0.6926 & 0.6027 \\
\texttt{no\_action} & 0.2059 & 0.2542 & 0.6089 & 0.3381 & 0.2270 & 0.1984 \\
\texttt{small\_maneuver} & 0.7467 & 0.7721 & 0.5000 & 0.9526 & 0.8604 & 0.6517 \\
\bottomrule
\end{tabular}
\caption{Valori medi delle feature per raccomandazione.}
\end{table}

La validazione mostra che:

\begin{itemize}
    \item \texttt{no\_action} ha il valore medio più basso di \texttt{risk\_score}.
    \item \texttt{monitor} è associata a rischio intermedio.
    \item \texttt{small\_maneuver} è associata a rischio alto e urgenza alta, ma priorità missione più bassa.
    \item \texttt{major\_maneuver} è associata a rischio alto e priorità missione più alta.
\end{itemize}

Questo conferma che il recommender è coerente e spiegabile.

% =========================
\section{Stato attuale del sistema}
% =========================

Allo stato attuale, il progetto implementa un sistema ibrido composto da due livelli principali:

\begin{itemize}
    \item un modulo di Machine Learning per la predizione di \texttt{max\_risk\_estimate};
    \item un recommendation layer rule-based e spiegabile per la selezione dell'azione finale.
\end{itemize}

La pipeline è completamente riproducibile tramite:

\begin{verbatim}
python3 main.py
\end{verbatim}

Eseguendo questo comando vengono generati il dataset completo, il file pulito delle raccomandazioni, il summary report, il validation report, il report di valutazione del modello ML, il modello addestrato e le visualizzazioni di esempio.

Il progetto è quindi completo come prototipo interpretabile di supporto decisionale con una componente predittiva basata su Machine Learning e una componente decisionale basata su regole spiegabili.

% =========================
\section{Limiti del progetto}
% =========================

La versione attuale del progetto presenta alcuni limiti:

\begin{itemize}
    \item il recommendation layer è rule-based e progettato manualmente, mentre il modulo di Machine Learning predice una stima di rischio del dataset e non decisioni reali degli operatori;
    \item \texttt{fuel\_cost} e \texttt{mission\_priority} sono assunzioni sintetiche;
    \item un utilizzo operativo reale richiederebbe soglie validate, costi reali di manovra, vincoli del satellite e revisione da parte di esperti;
    \item i pesi degli score sono progettati e calibrati manualmente;
    \item il sistema non simula manovre orbitali;
    \item il sistema non predice esiti reali di collisione, ma supporta la decisione;
    \item le visualizzazioni sono statiche e non rappresentano una propagazione orbitale completa.
\end{itemize}

% =========================
\section{Sviluppi futuri}
% =========================

Possibili sviluppi futuri includono:

\begin{itemize}
    \item migliorare il modulo di Machine Learning confrontando modelli aggiuntivi e ottimizzando gli iperparametri;
    \item stimare il costo di manovra con calcoli reali di delta-v;
    \item usare metadati reali sulla priorità della missione;
    \item validare le regole di scoring con decisioni etichettate da esperti;
    \item aggiungere una dashboard interattiva;
    \item creare animazioni temporali usando più osservazioni dello stesso evento;
    \item includere analisi dell'incertezza tramite feature di covarianza;
    \item confrontare le raccomandazioni rule-based con decisioni storiche degli operatori.
\end{itemize}

% =========================
\section{Conclusione}
% =========================

Questo progetto costruisce con successo un sistema di raccomandazione spiegabile per il supporto decisionale nella collision avoidance satellitare.

Il sistema trasforma dati grezzi di conjunction in raccomandazioni a livello evento. Crea feature fisiche e operative, calcola un risk score, valuta le azioni candidate e seleziona una raccomandazione finale.

Il progetto include inoltre un modulo di Machine Learning per la predizione di \texttt{max\_risk\_estimate}. Il modello predittivo utilizza feature fisiche e orbitali dello scenario per stimare il rischio massimo dell'evento. Il modello non predice direttamente l'azione finale, perché il dataset non contiene decisioni reali degli operatori come etichette.

La raccomandazione finale continua a essere prodotta dal recommendation layer spiegabile, che combina rischio, distanza, urgenza, velocità, costo sintetico e priorità missione per selezionare l'azione più adatta.

Il prototipo attuale produce output tabellari, report, un modello ML salvato e visualizzazioni di esempio di scenari rappresentativi. Può essere considerato una prima versione completa e interpretabile di un sistema ibrido per la collision avoidance satellitare, composto da una componente predittiva basata su Machine Learning e una componente decisionale rule-based.

\end{document}